% ---------------------------------------------------------
% Lean mapping note:
% HasLeastUpperBoundProperty, AxiomOfCompleteness,
% CompletenessGivesRealSupremum.
\subsection{Axiom of Completeness}
% ---------------------------------------------------------


\begin{remark*}[Interpretation]
The ordered field axioms alone do not prevent ``holes'' in the number line.
The rationals $\mathbb{Q}$ satisfy all field and order axioms, yet fail
completeness. Consider the set
\[
S = \{x \in \mathbb{Q} : x^2 < 2\}.
\]
This set is nonempty (e.g.\ $1 \in S$) and bounded above in $\mathbb{Q}$
(e.g.\ $2$ is an upper bound). Yet $\sup S$ does not exist as a rational
number: the candidate $\sqrt{2}$ is irrational. The set $S$ has no least
upper bound in $\mathbb{Q}$ --- a hole sits exactly where $\sqrt{2}$ should be.

The Completeness Axiom asserts that $\mathbb{R}$ has no such holes:
every nonempty bounded-above subset of $\mathbb{R}$ has a supremum
\emph{inside} $\mathbb{R}$.  This single axiom is what distinguishes
$\mathbb{R}$ from $\mathbb{Q}$, and it underlies every major theorem in
analysis.
\end{remark*}

\begin{remark*}[Lean alignment]
In the Lean formalization for this chapter, the analysis-facing declaration is
\[
  \mathsf{AxiomOfCompleteness} :
  \mathsf{HasLeastUpperBoundProperty}\ \mathbb{R}.
\]
Thus Volume III assumes completeness of the real line directly as a structural
property of $\mathbb{R}$.
\end{remark*}


\begin{definitionbox}{Definition (Least Upper Bound Property)}
\begin{definition}[Least Upper Bound Property]
\label{def:least-upper-bound-property}
Let $(S,\le)$ be a partially ordered set.  The set $S$ has the
\emph{least upper bound property} if every nonempty subset of $S$ that is
bounded above in $S$ has a supremum in $S$.
\end{definition}
\end{definitionbox}
\LeanFormalizes{def:least-upper-bound-property}{lra-lean}{LRA.VolumeIII.Analysis.Completeness.AxiomOfCompleteness}{HasLeastUpperBoundProperty}{pending}

\begin{remark*}[Standard quantified statement]
\[
\begin{aligned}
&\operatorname{HasLeastUpperBoundProperty}(P)
\Longleftrightarrow
\forall A\subseteq S\;
\Bigl[
\bigl(A\neq\varnothing \land \exists u\in S\;\forall x\in A\;(x\le u)\bigr)
\Longrightarrow
\exists s\in S\;
\Bigl[
(\forall a\in A)(a\leq s)
\land
\forall u\in S\;
\bigl((\forall a\in A)(a\leq u)\Rightarrow s\leq u\bigr)
\Bigr]
\Bigr].
\end{aligned}
\]
\end{remark*}

\begin{remark*}[Predicate reading]
\[
S\in\mathbf{Set},\qquad P=\mathsf{PartiallyOrderedSet}(S,\leq),\qquad
\operatorname{HasLeastUpperBoundProperty}(P)\coloneqq
\forall A\subseteq S\;
\Bigl[
\bigl(A\neq\varnothing \land \exists u\in S\;\forall x\in A\;(x\le u)\bigr)
\rightarrow
\exists s\in S\;\operatorname{LeastUpperBound}(s,A,P)
\Bigr].
\]
\end{remark*}

\begin{remark*}[Negated quantified statement]
\[
\begin{aligned}
&\text{Let }(S,\le)\text{ be a partially ordered set}.\\
&S\text{ does not have the Least Upper Bound Property}\\
&\Longleftrightarrow
\exists A\subseteq S\;
\Bigl[
A\neq\varnothing
\land
\exists u\in S\;\forall x\in A\;(x\le u)
\land
\forall s\in S\;\neg\operatorname{LeastUpperBound}(s,A)
\Bigr].
\end{aligned}
\]
\end{remark*}

\begin{remark*}[Negation predicate reading]
\[
S\in\mathbf{Set},\qquad P=\mathsf{PartiallyOrderedSet}(S,\leq),\qquad
\neg\operatorname{HasLeastUpperBoundProperty}(P)\Longleftrightarrow
\exists A\subseteq S\;
\Bigl[
A\neq\varnothing
\land
\exists u\in S\;\forall x\in A\;(x\le u)
\land
\forall s\in S\;\neg\operatorname{LeastUpperBound}(s,A,P)
\Bigr].
\]
\end{remark*}

\begin{remark*}[Failure modes]
\begin{description}
  \item[Exposition.]
\[
\neg\operatorname{HasLeastUpperBoundProperty}(S)
=
\underbrace{\exists A\subseteq S\;(A\neq\varnothing)}_{\text{nonempty witness set}}
\;\land\;
\underbrace{\exists u\in S\;\forall x\in A\;(x\le u)}_{\text{bounded above in }S}
\;\land\;
\underbrace{\forall s\in S\;\neg\operatorname{LeastUpperBound}(s,A)}_{\text{no supremum in }S}.
\]


  \item[\(\operatorname{LeastUpperBound} / \operatorname{HasLeastUpperBoundProperty}.\)]
  \[
\begin{aligned}
&\text{Let }(S,\le)\text{ be a partially ordered set}.\\
&S\text{ does not have the Least Upper Bound Property}\\
&\Longleftrightarrow
\exists A\subseteq S\;
\Bigl[
A\neq\varnothing
\land
\exists u\in S\;\forall x\in A\;(x\le u)
\land
\forall s\in S\;\neg\operatorname{LeastUpperBound}(s,A)
\Bigr].
\end{aligned}
\]
\[
S\in\mathbf{Set},\qquad P=\mathsf{PartiallyOrderedSet}(S,\leq),\qquad
\neg\operatorname{HasLeastUpperBoundProperty}(P)\Longleftrightarrow
\exists A\subseteq S\;
\Bigl[
A\neq\varnothing
\land
\exists u\in S\;\forall x\in A\;(x\le u)
\land
\forall s\in S\;\neg\operatorname{LeastUpperBound}(s,A,P)
\Bigr].
\]
\end{description}
\end{remark*}

\begin{remark*}[Interpretation]
The Least Upper Bound Property says that upper-bounded nonempty subsets do not point toward missing ceiling elements. Whenever the order supplies at least one upper bound, it also supplies a smallest upper bound in the same ambient set. Failure therefore indicates a hole in the ordered structure rather than a defect in the subset alone.
\end{remark*}

\begin{dependencies}
\begin{itemize}
  \item \hyperref[def:partially-ordered-set]{Partially Ordered Set}
  \item \hyperref[def:subset]{Subset}
  \item \hyperref[def:real-upper-bound]{Upper Bound}
  \item \hyperref[def:supremum]{Supremum}
\end{itemize}
\end{dependencies}






\begin{axiombox}{Axiom (Completeness of $\mathbb{R}$)}
\begin{axiom}[Completeness of $\mathbb{R}$]
\label{ax:completeness-of-reals}
Every nonempty subset $A$ of $\mathbb{R}$ that is bounded above has a
supremum in $\mathbb{R}$.
\end{axiom}
\end{axiombox}

\begin{remark*}[Standard quantified statement]
\[
\operatorname{HasLeastUpperBoundProperty}(\mathbb{R}).
\]
Equivalently,
\[
\begin{aligned}
&\forall A\subseteq\mathbb{R}\;
\Bigl[
\bigl(A\neq\varnothing \land \exists u\in\mathbb{R}\;\forall x\in A\;(x\le u)\bigr)
\rightarrow
\exists s\in\mathbb{R}\;\operatorname{LeastUpperBound}(s,A)
\Bigr].
\end{aligned}
\]
\end{remark*}

\begin{remark*}[Predicate reading]
\[
P_{\mathbb{R}}=\mathsf{OrderedSet}(\mathbb{R},\leq),\qquad
\operatorname{HasLeastUpperBoundProperty}(P_{\mathbb{R}}).
\]
\end{remark*}

\begin{remark*}[Negated quantified statement]
\[
\begin{aligned}
&\neg\operatorname{HasLeastUpperBoundProperty}(\mathbb{R})\\
&\Longleftrightarrow
\exists A\subseteq\mathbb{R}\;
\Bigl[
A\neq\varnothing
\land
\exists u\in\mathbb{R}\;\forall x\in A\;(x\le u)
\land
\forall s\in\mathbb{R}\;\neg\operatorname{LeastUpperBound}(s,A)
\Bigr].
\end{aligned}
\]
\end{remark*}

\begin{remark*}[Negation predicate reading]
\[
P_{\mathbb{R}}=\mathsf{OrderedSet}(\mathbb{R},\leq),\qquad
\neg\operatorname{HasLeastUpperBoundProperty}(P_{\mathbb{R}})\Longleftrightarrow
\exists A\subseteq\mathbb{R}\;
\Bigl[
A\neq\varnothing
\land
\exists u\in\mathbb{R}\;\forall x\in A\;(x\le u)
\land
\forall s\in\mathbb{R}\;\neg\operatorname{LeastUpperBound}(s,A,P_{\mathbb{R}})
\Bigr].
\]
\end{remark*}

\begin{remark*}[Failure modes]
\begin{description}
  \item[Exposition.]
\[
\neg\operatorname{HasLeastUpperBoundProperty}(\mathbb{R})
=
\underbrace{\exists A\subseteq\mathbb{R}\;(A\neq\varnothing)}_{\text{nonempty witness set}}
\;\land\;
\underbrace{\exists u\in\mathbb{R}\;\forall x\in A\;(x\le u)}_{\text{bounded above in }\mathbb{R}}
\;\land\;
\underbrace{\forall s\in\mathbb{R}\;\neg\operatorname{LeastUpperBound}(s,A)}_{\text{missing real supremum}}.
\]


  \item[\(\operatorname{HasLeastUpperBoundProperty} / \operatorname{LeastUpperBound}.\)]
  \[
\begin{aligned}
&\neg\operatorname{HasLeastUpperBoundProperty}(\mathbb{R})\\
&\Longleftrightarrow
\exists A\subseteq\mathbb{R}\;
\Bigl[
A\neq\varnothing
\land
\exists u\in\mathbb{R}\;\forall x\in A\;(x\le u)
\land
\forall s\in\mathbb{R}\;\neg\operatorname{LeastUpperBound}(s,A)
\Bigr].
\end{aligned}
\]
\[
P_{\mathbb{R}}=\mathsf{OrderedSet}(\mathbb{R},\leq),\qquad
\neg\operatorname{HasLeastUpperBoundProperty}(P_{\mathbb{R}})\Longleftrightarrow
\exists A\subseteq\mathbb{R}\;
\Bigl[
A\neq\varnothing
\land
\exists u\in\mathbb{R}\;\forall x\in A\;(x\le u)
\land
\forall s\in\mathbb{R}\;\neg\operatorname{LeastUpperBound}(s,A,P_{\mathbb{R}})
\Bigr].
\]
\end{description}
\end{remark*}

\begin{remark*}[Interpretation]
The completeness axiom asserts that the real line has no order gaps of the
least-upper-bound kind. Nonempty sets of real numbers that are bounded above
always determine a real number sitting exactly at the least possible ceiling.
It is called an axiom here because, in this analysis volume, it is assumed as
a structural property of the real line. It is not proved from the ordered-field
axioms alone. Volume II may construct real-number models and prove analogous
completeness properties for those models; Volume III uses completeness of
$\mathbb{R}$ directly as the standing least-upper-bound principle for real
analysis. This is the structural property that distinguishes $\mathbb{R}$ from
ordered fields such as $\mathbb{Q}$.
\end{remark*}

\begin{dependencies}
\begin{itemize}
  \item \hyperref[def:reals]{The Real Numbers}
  \item \hyperref[def:real-ordered-field]{Ordered Field}
  \item \hyperref[def:ordered-set]{Ordered Set}
  \item \hyperref[def:subset]{Subset}
  \item \hyperref[def:least-upper-bound-property]{Least Upper Bound Property}
\end{itemize}
\end{dependencies}

% Lean-aligned additions inserted for Volume III Lean-to-TeX handoff.
% Lean-aligned addition: thm:exists-unique-nonneg-sqrt
\begin{theorem}[Existence and Uniqueness of Nonnegative Square Roots]
\label{thm:exists-unique-nonneg-sqrt}
For every $a\in\mathbb R$ with $a\ge0$, there is a unique $b\in\mathbb R$ such that $b\ge0$ and $b^2=a$.

\smallskip
\noindent
\hyperref[prf:exists-unique-nonneg-sqrt]{\textit{Go to proof.}}
\end{theorem}
\LeanFormalizes{thm:exists-unique-nonneg-sqrt}{lra-lean}{LRA.VolumeIII.Analysis.Completeness.CompletenessAdditions}{ExistsUniqueNonnegSqrt}{pending}

\begin{remark*}[Standard quantified statement]
For every $a\in\mathbb R$ with $a\ge0$, there is a unique $b\in\mathbb R$ such that $b\ge0$ and $b^2=a$.
\end{remark*}

\begin{remark*}[Interpretation]
This formal object records a Lean-aligned addition in the local language of this section.
\end{remark*}

\begin{dependencies}
\NoLocalDependencies
\end{dependencies}
% Lean-aligned addition: thm:rationals-lack-lub-property
\begin{theorem}[Rationals Lack the Least Upper Bound Property]
\label{thm:rationals-lack-lub-property}
The ordered field $\mathbb Q$ does not have the least upper bound property. In particular, the set $\{q\in\mathbb Q:q^2<2\}$ is nonempty and bounded above in $\mathbb Q$, but it has no least upper bound in $\mathbb Q$.

\smallskip
\noindent
\hyperref[prf:rationals-lack-lub-property]{\textit{Go to proof.}}
\end{theorem}
\LeanFormalizes{thm:rationals-lack-lub-property}{lra-lean}{LRA.VolumeIII.Analysis.Completeness.CompletenessAdditions}{RationalsLackLubProperty}{pending}

\begin{remark*}[Standard quantified statement]
The ordered field $\mathbb Q$ does not have the least upper bound property. In particular, the set $\{q\in\mathbb Q:q^2<2\}$ is nonempty and bounded above in $\mathbb Q$, but it has no least upper bound in $\mathbb Q$.
\end{remark*}

\begin{remark*}[Interpretation]
This formal object records a Lean-aligned addition in the local language of this section.
\end{remark*}

\begin{dependencies}
\NoLocalDependencies
\end{dependencies}
